\documentclass[10pt,a4paper]{article}
\usepackage[margin=22mm,headheight=15pt]{geometry}
\usepackage{fontspec}
\setmainfont{Latin Modern Roman}
\setsansfont{DejaVu Sans}
\setmonofont{DejaVu Sans Mono}[Scale=0.80]
\usepackage{amsmath,amssymb,amsthm,booktabs,tabularx,array,enumitem,xcolor,fancyhdr,hyperref,ragged2e}
\definecolor{ink}{HTML}{112C4D}
\definecolor{blue}{HTML}{176C9A}
\hypersetup{colorlinks=true,urlcolor=blue,linkcolor=blue,pdftitle={FT1536: rozpoznanie QROM, S05-Q001},pdfauthor={Notatka robocza dla Niirmata}}
\pagestyle{fancy}\fancyhf{}
\fancyhead[L]{\sffamily\small FREE FT / S05}
\fancyhead[R]{\sffamily\small QROM - ROZPOZNANIE 001}
\fancyfoot[L]{\sffamily\footnotesize Snapshot 2d2cdf5 / 10 października 2026}
\fancyfoot[R]{\sffamily\footnotesize\thepage}
\setlength{\parindent}{0pt}\setlength{\parskip}{6pt}
\setlength{\emergencystretch}{3em}
\setlist[itemize]{leftmargin=*,itemsep=3pt,topsep=3pt}
\setlist[enumerate]{leftmargin=*,itemsep=3pt,topsep=3pt}
\renewcommand{\arraystretch}{1.25}
\newcommand{\code}[1]{\texttt{\detokenize{#1}}}
\newcommand{\head}[1]{\section*{\color{ink}\sffamily #1}}
\newcommand{\subhead}[1]{\subsection*{\color{ink}\sffamily #1}}
\newcommand{\E}{\mathbb E}
\newcommand{\Tr}{\operatorname{Tr}}
\newcommand{\PhiF}{\Phi}
\newcommand{\base}{https://github.com/niirmataa/free_falcon_sign/blob/2d2cdf52304cee9cf513487d4296392262df22f4/}
\begin{document}
\begin{center}
{\sffamily\color{blue}\small NOTATKA BADAWCZA / CZĘŚĆ II / S05}\par\vspace{8pt}
{\sffamily\bfseries\color{ink}\LARGE FT1536 w QROM}\par\vspace{5pt}
{\sffamily\large Droga przez pełny rozkład odpowiedzi Sign}\par\vspace{8pt}
{\small Dla Niirmata / stan rozpoznania: 10 października 2026}
\end{center}
\head{1. Wynik tego rozpoznania}
\textbf{Istnieje konkretny kierunek dla niezmienionego, ograniczonego do 16 prób Sign: porównywać całe odpowiedzi przy jednym nonce.} Nie wykazano jeszcze jego bezpieczeństwa w QROM. Otrzymano mapę rzeczywistych interfejsów oraz elementarny lemat: kierunkowa nierówność z $\Phi$ pozostaje ważna po wspólnej kwantowej kontynuacji klasycznego losowania. Pełny dowód tekstowy jest na stronie~4.

Trzy ustalenia zmieniają punkt startowy:
\begin{enumerate}
\item \textbf{Ekstrakcja w obecnym modelu nie używa rewindingu.} Konstruktor \code{Reduction.build} wykonuje symulację i stosuje \code{finishSim}. Świadek wynika z pojedynczego zaakceptowanego fałszerstwa. Cztery strzałki paperu są przejściami między eksperymentami.
\item \textbf{Pierwszy cel to kwantowy hash i klasyczny Sign.} Spójny dostęp do Sign jest innym, silniejszym modelem. Nie jest wymagany do tego celu QROM [2].
\item \textbf{Nonce oraz hash pozostają stałe podczas prób.} Wobec tego gotowego twierdzenia dla CoreFalcon$^+$ nie można zastosować przez samo podstawienie parametrów. Rysunek~9 w [1] pokazuje różnicę algorytmów.
\end{enumerate}

\subhead{Zakres i status}
Rozpoznanie rozwija istniejący wpis \textbf{S05 - QROM} w \code{docs/onboarding/ROADMAP.md}. Nie wznawia B1.07 ani pozostałych aktywnych pakietów. Nie zmienia C, Lean, paperu, strony WWW ani stanu repozytorium. Dalsza integracja jest osobną decyzją; pierwszeństwo głównej linii ePrint pozostaje zachowane.

\textbf{Baza:} \code{niirmataa/free_falcon_sign}, commit
\begin{center}\code{2d2cdf52304cee9cf513487d4296392262df22f4}.\end{center}
Przeczytano wybrane źródła i instrukcje; 23 lokalne kopie źródeł sprawdzono z identyfikatorami blobów Git. Jest to kontrola integralności odczytu, nie replay dowodów. Lean i Sage nie były dostępne w tym środowisku. Nowy lemat ma wyłącznie dowód tekstowy i wewnętrzne sprawdzenie, bez statusu kernelowego ani projektowego REVIEWED.

\newpage
\head{2. Co naprawdę robi przypięty Sign}
Odczyt \code{Extra/c/falcon-sign.c} daje następującą kolejność:
\begin{enumerate}
\item \code{falcon_sign_start}: pobranie 40 bajtów nonce z generatora (3281--3288).
\item \code{falcon_sign_update}: dołączenie wiadomości do kontekstu hasha (3301--3303).
\item \code{falcon_sign_generate}: pojedyncze hash-to-point do \code{hm}, przed pętlą (3324--3325).
\item Maksymalnie 16 wywołań \code{do_sign}, zawsze z tym samym \code{hm} (cap 101--109; pętla 3330--3409).
\item Błąd samplera kończy wywołanie (3374--3375). Odrzucenie normy ponawia próbę. Po akceptacji następuje jedno kodowanie; jego porażka jest końcowa (3411--3421).
\end{enumerate}
\textbf{Wniosek źródłowy:} liczba prób samplera, liczba nonce i liczba obliczeń hasha nie są tym samym licznikiem. Domyślny interfejs losuje nonce raz na podpis. Dostępne jest także API z zewnętrznym nonce; własność entropii nie obejmuje dowolnego wejścia do tego API.

\subhead{Obiekt porównania}
Dla ustalonego klucza publicznego $h$ i wyzwania $c$ proponujemy zachować cały rozkład
\[
 P_h(c,o)=\frac{1}{|R_q|}\,\operatorname{signBody}_h(c)(o),
 \qquad o\in\operatorname{Option}(\operatorname{BoxVec}).
\]
To dokładna postać \code{honest_joint_body} z \code{Run2/LawBinding.lean}. Wartość \code{none} obejmuje porażkę modelowego ciała Sign. Nie wolno warunkować całego eksperymentu na wydaniu podpisu.

Dla modelowego prawa jednej próby $T_c$, $u_c=T_c(\mathrm{none})$ i mapy emisji $E$, eksport \code{signBodyOf_mass_none} daje
\[
 \Pr[\mathrm{none}\mid c]=u_c^{16}
 +\left(\sum_{i=0}^{15}u_c^i\right)
   \sum_{z:E(z)=\mathrm{none}}T_c(\mathrm{some}\ z).
\]
Pierwszy składnik to wyczerpanie prób, drugi to końcowa porażka emisji. Tożsamość dotyczy modelowego prawa, nie stanowi sama dowodu zgodności z terminalnym błędem samplera, błędem bufora i wszystkimi zwrotami C. Te zachowania wymagają osobnej realizacji źródłowej.

\subhead{Ekstrakcja bez cofania wykonania}
\code{Relation.extract} definiuje
\[
 z^*=\bigl(\operatorname{center}(c^*-h\operatorname{reduce}(s^*)),\ s^*\bigr).
\]
\code{accepted_extracts} dowodzi: $\operatorname{Verify}(h,c^*,s^*)\Rightarrow
\operatorname{ShortPreimage}(h,c^*,z^*)$. Równanie zachowuje sens dla klasycznego wyniku kwantowego przeciwnika. Otwartym krokiem QROM jest osadzenie właściwego celu i kontrola sukcesu redukcji; klasyczny indeks z tablicy ROM nie przenosi się automatycznie.

\newpage
\head{3. Mapa lematów: zachować czy rozszerzyć?}
Prefiks \code{R/} oznacza \code{proofs/ft1536/development/T12_1/run2/formal/}.

\begin{tabularx}{\linewidth}{@{}>{\raggedright\arraybackslash}p{.34\linewidth}>{\raggedright\arraybackslash}X@{}}
\toprule
\textbf{Element} & \textbf{Ocena dla S05}\\\midrule
\code{Relation.lean}: \code{extraction_equation}, \code{accepted_extracts}
& Zachować twierdzenia algebraiczne. Nie dostarczają kwantowej symulacji ani osadzenia celu.\\
\code{R/SignLayerSupport.lean}: masy \code{some}/\code{none}, \code{cap_le_of_pointwise}
& Zachować jako rachunek klasycznego, prywatnego losowania Sign. Sprawdzić identyczność prawa i granicy prób.\\
\code{R/AttemptPointwise.lean}: \code{AttemptShape}, \code{Sandwich}
& Zachować interfejsy jako warunkowe przesłanki. Nie uznawać realizacji źródłowej za ukończoną.\\
\code{R/Run2/LawBinding.lean}: \code{LocalJointCertificate}
& Obejmuje wspólne prawo wyzwania i odpowiedzi, łącznie z porażką. Potrzebuje interpretacji w nowym eksperymencie.\\
\code{EventTransfer.lean}: \code{event_quadratic}, \code{phi_bound}
& Algebra pozostaje użyteczna. Lemat na stronie~4 rozszerza test wskaźnikowy do prawdopodobieństwa akceptacji $v(x)\in[0,1]$.\\
\code{R/Run2/Games.lean}: \code{Program}, \code{ClassicalAdversary}
& Nowa semantyka kwantowa konieczna. Obecny typ pozwala wyłącznie na klasyczne zapytania i gałęzie.\\
\code{R/Run2/StoppedComparison.lean}: \code{programComparison}
& Dowód jest indukcją po klasycznym \code{Program}. Nie daje bez dodatkowej pracy potęgi $(1+e)^{q_s}$ w QROM.\\
Klasyczna świeżość, \code{hashTargets}, $q_h+1$ celów
& Zastąpić dowodem reprogramowania i osadzenia wyzwania. Nie kopiować straty kolizyjnej ani liczby celów.\\
Hybryda rzeczywistego PRG
& Osobny kontrakt przeciw kwantowemu rozróżniaczowi; obejmuje rzeczywisty horyzont taśmy i te same wyniki porażki.\\
\bottomrule
\end{tabularx}

\subhead{Dodatkowa granica: dostęp do stanu}
\code{Sampler.code} przyjmuje klasyczny \code{State}, zawierający tablicę ROM i zdarzenia. W QROM nie wolno odczytywać wszystkich zapytań przeciwnika w celu zbudowania takiego stanu. Potrzebny jest efektywny sampler zależny tylko od dozwolonej informacji albo nowy, wykonalny interfejs kwantowej redukcji.

\textbf{Kierunek momentu jest ustalony:}
\[
 J=\operatorname{samplerLaw}S,\quad P=\operatorname{freshHonest}h,\qquad
 J\ll P,\quad\sum_x\frac{J(x)^2}{P(x)}\le 1+e.
\]
Ten kierunek pozwala ograniczać sukces w $P$ przez sukces w $J$ za pomocą $\Phi$. Odwrócenie ilorazu zmienia wniosek.

\newpage
\head{4. Pierwszy lemat: wspólna kontynuacja kwantowa}
\textbf{Status: pełny argument tekstowy; elementarny lemat pomocniczy, bez roszczenia nowości i bez formalizacji Lean.}

Niech $X$ będzie skończone, $P,J$ znormalizowanymi rozkładami, $J\ll P$, $D\ge0$ oraz
\[
 \sum_{x:P(x)>0}\frac{J(x)^2}{P(x)}\le 1+D.
\]
Dla każdego $x$ niech $\tau_x$ będzie znormalizowanym stanem kwantowym w tej samej skończenie wymiarowej przestrzeni. \textbf{Rodzina $\tau_x$ jest wspólna dla obu eksperymentów.} Może obejmować cały wynik kontynuacji, łącznie z dozwoloną wyrocznią. Niech $0\le M\le I$ będzie wspólnym efektem pomiarowym. Zdefiniujmy
\[
 \rho_P=\sum_xP(x)\tau_x,\quad \rho_J=\sum_xJ(x)\tau_x,
 \qquad a=\Tr(M\rho_P),\quad b=\Tr(M\rho_J).
\]
Wtedy
\[
 (a-b)^2\le D\,a(1-a),\qquad
 a\le\PhiF(D,b)=\frac{2b+D+\sqrt{D^2+4Db(1-b)}}{2(1+D)}.
\]

\textbf{Dowód.} Niech $v(x)=\Tr(M\tau_x)\in[0,1]$ oraz $L(x)=J(x)/P(x)$ na nośniku $P$. Mamy $\E_P L=1$, $a=\E_Pv$ i $b=\E_P(Lv)$. Stąd
\[
 b-a=\E_P[(L-1)(v-a)],\qquad
 \E_P[(L-1)^2]=\E_P[L^2]-1\le D.
\]
Ponieważ $v^2\le v$,
\[
 \E_P[(v-a)^2]=\E_P[v^2]-a^2\le a-a^2.
\]
Cauchy--Schwarz daje nierówność kwadratową. Po uporządkowaniu:
\[
 (1+D)a^2-(2b+D)a+b^2\le0.
\]
Górny pierwiastek to podane $\PhiF(D,b)$. Dla $D=0$ otrzymujemy $a=b$; zera masy nie stanowią problemu dzięki $J\ll P$.\hfill$\square$

\subhead{Dlaczego to pomaga i czego jeszcze nie daje}
Kwantowa kontynuacja sama w sobie nie niszczy tego rachunku. Trzeba jednak \emph{skonstruować} wspólne $\tau_x$ dla konkretnych gier oraz wykazać kontrolę całego rozkładu $X$. Lemat nie dowodzi tej reprezentacji, adaptacyjnego składania ani efektywności samplera.

\textbf{Kontrola konieczności przesłanki.} Niech $b$ będzie uczciwym bitem w obu grach. W pierwszej przekażmy odbiorcy $(b,b)$, w drugiej $(b,1-b)$. Marginały pierwszej składowej są identyczne ($D=0$), ale test równości pary ma sukces $1$ i $0$. Już klasyczna informacja dodatkowa wystarcza, by obalić wnioskowanie tylko z marginału. To przykład diagnostyczny, nie atak na FT1536 ani nowy wynik negatywny.

\newpage
\head{5. Proponowana konstrukcja eksperymentu}
\textbf{Cel początkowy:} EUF-CMA dla modelu współczynnikowego FT1536, z kwantowym dostępem do direct-output hasha $H:\mathcal X\to R_q$, klasycznym interfejsem Sign oraz klasycznym fałszerstwem $(m^*,r^*,s^*)$. Długości wejść i zasoby są jawnie ograniczone. Ustalona kodowana reprezentacja $R_q$ określa działanie wyroczni unitarnej. Publiczny SHAKE/H2P pozostaje osobnym torem S04.

\begin{enumerate}
\item \textbf{Gra Q0:} rzeczywista dla wybranego modelu matematycznego. Jeden świeży idealny nonce na zlecenie, jedno wyzwanie, maksymalnie 16 prywatnych prób, jawny wynik porażki. Wiadomość trafia do zbioru zapytań Sign także przy porażce.
\item \textbf{Gra Q1:} przeprogramować punkt $(r,m)$ na świeże jednostajne wyzwanie, zachowując całe warunkowe prawo odpowiedzi. Osobno wykazać stratę tego przejścia wobec kwantowego przeciwnika.
\item \textbf{Gra Q2:} zastąpić joint law $(c,o)$ prawem publicznej symulacji, korzystając z właściwie skierowanego momentu. Cała porażka należy do $o$. Konstrukcja musi zapewnić wspólną kontynuację ze strony~4.
\item \textbf{Gra Q3:} osadzić niezależny cel ISIS/MT-ISIS i wyprowadzić świadek z końcowego wyniku, z nowym rachunkiem czasu, liczby zapytań i straty sukcesu.
\end{enumerate}
To \textbf{szkic gier}, nie ukończony łańcuch redukcji. Ich wspólna semantyka i zgodność z realnym Sign wymagają dowodu.

\subhead{Cztery warunki rozstrzygające drogę}
\begin{itemize}
\item \textbf{Q-JOINT:} efektywne, bezpułapkowe losowanie $(c,o)$ o właściwej masie każdej porażki. Zgodność z abstract \code{Sampler} nie wystarcza bez kosztu i legalnego dostępu do stanu.
\item \textbf{Q-COMMON:} wspólna kontynuacja po tym samym klasycznym wyniku, także przy kwantowym stanie przeciwnika i wyroczni. Nie wystarcza równość marginałów.
\item \textbf{Q-COMPOSE:} adaptacyjny rachunek momentu dla tego interfejsu. Dopiero jego dowód może uzasadnić $D_Q\le(1+e)^{q_s}-1$.
\item \textbf{Q-TARGET:} pełna konstrukcja reduktora i właściwa dystrybucja jego klucza/celów. Sam \code{accepted_extracts} nie ustala tych rozkładów.
\end{itemize}

\textbf{Liczniki:} $q_s$ - zlecenia Sign; $C_s\le16q_s$ - próby samplera; $R$ - faktyczne reprogramowania; $Q$ - wszystkie zapytania uwzględnione w odpowiednim lemacie. Dla proponowanej gry chcemy $R=q_s$, lecz jest to wymaganie konstrukcji, a nie ustalony fakt. W idealnym wariancie świeżego nonce $p_{\max}=2^{-320}$; realny generator wymaga osobnej hybrydy.

\newpage
\head{6. Literatura i trzy sprawdzone kierunki}
\subhead{A. Gotowe twierdzenie dla podobnego schematu}
W [1], wersja~3, lemat~6 dotyczy celowanego przeprogramowania. Twierdzenie~5 obejmuje CoreFalcon$^+$: nowe nonce i wyzwanie w każdej próbie; liczy wszystkie próby. Rysunek~9 (s.~70) odróżnia ten wariant od stałego nonce. \textbf{Wynik sprawdzenia: brak identyczności algorytmu z FT1536.} Nie jest to dowód niemożliwości naszej drogi.

\subhead{B. Pełna odpowiedź i jedno reprogramowanie}
To kierunek wybrany do dalszego S05. Zachowuje pętlę FT1536 i korzysta z istniejącego joint certificate. Własny lemat ze strony~4 jest pozytywnym wynikiem pomocniczym. Pierwszym nierozwiązanym krokiem pozostaje Q-COMMON wraz z Q-JOINT. Do adaptacyjnego przeprogramowania punktem odniesienia jest [2], twierdzenie~1 i propozycja~2; potrzebna jest entropia w odpowiednim momencie eksperymentu.

\subhead{C. Ogólna zamiana całej wyroczni}
Twierdzenie~4 w [1] wyklucza uniwersalne, czysto mnożnikowe przeniesienie kontroli Rényiego na kwantowy dostęp do całej wyroczni. \textbf{Ten ogólny wynik negatywny jest znany.} Nowy kontrprzykład musiałby dotyczyć innej, precyzyjnej przesłanki FT1536. Nie ustalono nowości takiego wyniku.

\subhead{Rola wskazanych przez właściciela narzędzi}
\begin{itemize}
\item \textbf{Kiltz--Lyubashevsky--Schaffner [3]:} twierdzenia 3.3--3.4 dotyczą określonych protokołów identyfikacyjnych i transformacji Fiat--Shamira. Przed zastosowaniem trzeba dostarczyć odpowiedni protokół i jego przesłanki, w tym właściwe własności symulacji i entropii. Sam nonce FT1536 tego nie zapewnia.
\item \textbf{Don--Fehr--Majenz--Schaffner [4] oraz Don--Fehr--Majenz [5]:} rozróżniamy pracę CRYPTO~2019 od techniki 2.0 z CRYPTO~2020. Twierdzenie~2 w sprawdzonej wersji [5] daje narzędzie measure-and-reprogram ze stratą $(2q+1)^2$. To kandydat do osadzenia celu; nie buduje sam wyroczni Sign.
\item \textbf{Unruh [6]:} sprawdzony wczesny lemat O2H kontroluje rozróżnialność przez prawdopodobieństwo trafienia zapytania. Transformacja Unruha jest odrębnym narzędziem. Jej pełnego twierdzenia z pracy 2015 nie pozyskano w tym odczycie i nie użyto jako przesłanki.
\end{itemize}
\textbf{Nie ustalono:} nowości proponowanej drogi, pełnego QROM dla FT1536, parametrów bezpieczeństwa ani wykonalnej instancji wszystkich czterech kontraktów. Rozpoznanie nie jest audytem całych dowodów cytowanych prac.

\newpage
\head{7. Następny pakiet S05: dokładny kontrakt}
\textbf{Proponowana nazwa:} \code{S05_QROM_WHOLE_REPLY_001}. Jest to projekt zadania, bez zmiany oficjalnego rejestru ani uruchamiania istniejącego wykonawcy.

\textbf{Wejścia:} bieżący przypięty model Sign, \code{LocalJointCertificate}, ograniczenie prób 16, lemat ze strony~4 oraz wybrana, dokładnie przypięta wersja lematów reprogramowania.

\textbf{Główny cel:} skonstruować dwie gry o kwantowym dostępie do tej samej wyroczni i klasycznym Sign, różniące się tylko prawem pełnej odpowiedzi. Udowodnić reprezentację wspólnej kontynuacji i oszacowanie momentu dla adaptacyjnej kompozycji, bez czytania ukrytych zapytań kwantowych.

\textbf{Kryteria wyniku dodatniego:}
\begin{enumerate}
\item Jawne typy wiadomości, nonce, wyzwań, podpisów i wszystkich obserwowanych porażek; precyzyjna semantyka interfejsów.
\item Jeden nonce na podpis i to samo $c$ przez wszystkie próby. Nie zastępować algorytmu wariantem z odświeżaniem nonce.
\item Dozwolony sampler publiczny i policzony koszt. Każde użycie klasycznej tablicy ROM wymaga legalnej realizacji lub eliminacji.
\item Konstrukcja wspólnych warunkowych stanów/kanałów; lokalne ograniczenia jednolite po dozwolonej historii, kluczu i wiadomości.
\item Dowód dokładnego $D_Q$ i jasny zakres wniosku. Parametry reprogramowania oraz redukcja ISIS mogą pozostać osobnymi nazwanymi zobowiązaniami, bez deklaracji pełnego bezpieczeństwa.
\end{enumerate}

\textbf{Kryteria wyniku negatywnego:} jawne rozkłady, stan, wyrocznia i przeciwnik spełniające rzeczywiste przesłanki proponowanego lematu, wraz z naruszoną nierównością. Wskazać, czy obalono zbyt słaby interfejs, konkretną technikę, czy docelową tezę. Porównać z [1]; nie nazywać automatycznie nowym twierdzeniem istniejącego zakazu globalnego transferu.

\subhead{Co zachowuje pierwszeństwo głównej linii}
Brak zmian algorytmu dla dopasowania do literatury. Brak podmiany klasycznych eksportów. Brak awansu B1/B4/B5 na podstawie tej notatki. Najpierw osobny kontrakt i mały wynik, potem decyzja o formalizacji oraz integracji po uzgodnieniu z torem ePrint.

\textbf{Najbliższy rozstrzygający krok:} skonstruować Q-COMMON dla jednej odpowiedzi z nowym nonce, obejmującej parę $(c,o)$ i porażkę, przy dowolnym uprzednim kwantowym stanie wyroczni. Jeśli ten krok się powiedzie, przejść do adaptacyjnej kompozycji. Jeśli zawiedzie, zapisać minimalną brakującą przesłankę lub konkretny kontrprzykład.

\newpage
\head{8. Źródła i odtwarzalność odczytu}
\small
\textbf{Repozytorium.} Wszystkie odnośniki i blob SHA-1 znajdują się w dołączonym \code{source_manifest.json}. Bazą jest commit z pierwszej strony. Najważniejsze lokalizacje:
\begin{itemize}
\item \href{\base Extra/c/falcon-sign.c}{\code{Extra/c/falcon-sign.c}}: 101--109, 3281--3296, 3324--3421.
\item \href{\base proofs/ft1536/development/T12_1/run2/formal/Run2/Games.lean}{\code{R/Run2/Games.lean}}: typy przeciwnika i samplera, \code{finishSim}, \code{Reduction.build}.
\item \href{\base proofs/ft1536/development/T12_1/run2/formal/Run2/LawBinding.lean}{\code{R/Run2/LawBinding.lean}}: \code{LocalJointCertificate}, \code{honest_joint_body}.
\item \href{\base proofs/ft1536/development/T12_1/run2/formal/SignLayerSupport.lean}{\code{R/SignLayerSupport.lean}}: \code{signBodyOf_mass_none}, \code{ReplyShape}, \code{e2}.
\item \href{\base proofs/ft1536/development/T12_1/run2/formal/Run2/StoppedComparison.lean}{\code{R/Run2/StoppedComparison.lean}}: klasyczna indukcja \code{programComparison}.
\item \href{\base proofs/ft1536/paper/sec_04_main.tex}{\code{paper/sec_04_main.tex}}: cztery przejścia i rozróżnienie celu od twierdzenia warunkowego.
\end{itemize}

\textbf{Literatura.} Numery twierdzeń dotyczą wyłącznie wymienionych wersji. Autoryzacja źródła i odczyt wskazanych fragmentów nie oznaczają niezależnego audytu całej pracy.
\begin{enumerate}[label={[\arabic*]},itemsep=7pt]
\item M. Beckmann, C. Majenz, \emph{Quantum Oracle Distribution Switching and Applications to Falcon and Ring Signatures}. arXiv:2602.16268v3, 9 IX 2026. Lemat 6, §§3.2--3.3; twierdzenie 5, §4; rysunek 9 i dodatek E.1. \href{https://arxiv.org/abs/2602.16268v3}{Wersja}, \href{https://arxiv.org/pdf/2602.16268}{PDF odczytany jako v3 na stronie tytułowej}.
\item A. B. Grilo, K. Hövelmanns, A. Hülsing, C. Majenz, \emph{Tight adaptive reprogramming in the QROM}. arXiv:2010.15103v2, 30 X 2020; ASIACRYPT 2021. Twierdzenie 1, propozycja 2, model z klasycznym SIGN. \href{https://arxiv.org/abs/2010.15103v2}{Wersja}, \href{https://arxiv.org/pdf/2010.15103}{PDF}.
\item E. Kiltz, V. Lyubashevsky, C. Schaffner, \emph{A Concrete Treatment of Fiat-Shamir Signatures in the Quantum Random-Oracle Model}. Accepted manuscript, 20 II 2018; EUROCRYPT 2018. Twierdzenia 3.3--3.4. \href{https://pure.uva.nl/ws/files/45974858/916.pdf}{Kopia uniwersytecka}.
\item J. Don, S. Fehr, C. Majenz, C. Schaffner, \emph{Security of the Fiat-Shamir Transformation in the Quantum Random-Oracle Model}. CRYPTO 2019; arXiv:1902.07556v4, 27 VII 2020. Twierdzenia 2 i 22. \href{https://arxiv.org/pdf/1902.07556}{PDF}.
\item J. Don, S. Fehr, C. Majenz, \emph{The Measure-and-Reprogram Technique 2.0: Multi-Round Fiat-Shamir and More}. CRYPTO 2020; arXiv:2003.05207v3, 7 III 2022. Twierdzenie 2. \href{https://arxiv.org/pdf/2003.05207}{PDF}.
\item D. Unruh, \emph{Revocable quantum time-capsules}, wczesny manuskrypt autora, 16 VIII 2013, lemat 31, s.40. \href{https://kodu.ut.ee/~unruh/publications/qtc.pdf}{PDF}. Nie utożsamiamy numeracji z późniejszą publikacją JACM.
\end{enumerate}

\textbf{Granice sprawdzenia.} Brak nowego replayu Lean/Sage, kompilacji C i estymacji bezpieczeństwa. Nie przetestowano jeszcze Q-COMMON ani Q-COMPOSE w formalnym modelu kwantowym. Pakiet zawiera własny tekst, metadane źródeł i checkpoint; nie redystrybuuje cudzych PDF-ów ani całego repozytorium.
\end{document}
